% Section 5.1: Introduction
% Auto-generated from markdown drafts

\section{Introduction}
\label{sec:5.1}

The preceding chapter examined the Galactic Center Excess, the brightest expected signature of dark matter annihilation in gamma rays. That analysis confronted a difficulty common to all Galactic Center studies: disentangling a potential dark matter signal from the complex astrophysical emission of the inner Milky Way, where overlapping source populations and uncertain diffuse backgrounds prevent definitive conclusions. We now turn to a complementary search strategy, one that targets dark matter in a very different environment.

Rather than analyzing a diffuse excess embedded in the bright Galactic foreground, this chapter asks whether any of the point sources detected by the Fermi Large Area Telescope could be dark matter subhalos. The $\Lambda$CDM cosmological model predicts that the Milky Way's dark matter halo hosts a vast population of gravitationally bound substructures~\cite{Diemand:2008in, Springel:2008cc}. The great majority of these subhalos have masses well below the threshold for retaining baryonic matter~\cite{Zavala:2019gpd} and consequently produce no electromagnetic radiation through conventional astrophysical processes. If, however, the dark matter particle annihilates into Standard Model final states, the densest of these subhalos could generate detectable gamma-ray emission. They would then appear in Fermi-LAT catalogs as unassociated point sources lacking any multi-wavelength counterpart.

This chapter proceeds as follows. We first review the theoretical predictions for the Milky Way's subhalo population within $\Lambda$CDM (Section~\ref{sec:5.2}) and their expected gamma-ray properties (Section~\ref{sec:5.3}). We then introduce the unassociated source problem in the fourth Fermi-LAT source catalog and survey previous attempts to identify dark matter subhalo candidates through classification-based methods (Section~\ref{sec:5.4}). These approaches share a common limitation: they rely on supervised classifiers trained under the assumption that the labeled and unlabeled source populations follow identical distributions---an assumption we show to be violated. Section~\ref{sec:5.5} presents the methodological response to this challenge: a quantification learning framework that models class prevalences and distribution shifts simultaneously, yielding a well-defined likelihood suitable for statistical hypothesis testing. The remainder of the chapter presents the full analysis, following~\cite{Amerio:2025dmh}.
